\chapter{Bridges to Neighboring Frameworks}
\label{ch:bridges}

The world-model calculus does not exist in a vacuum.  It touches---and is
touched by---a constellation of neighboring formalisms: Markov Logic
Networks, probabilistic logic programming, probabilistic programming
languages, the classical PLN rule family,
NARS, first-order logic, categorical probability theory, and process
calculi.  This chapter constructs explicit bridges between each of these
and the world-model calculus.  A \emph{bridge} is not a metaphor.  It is a
pair of translation functions (embedding and projection) together with a
theorem stating what structural properties are preserved.

The goal is twofold: to show practitioners of each neighboring framework
exactly how their objects map into (or out of) the WM calculus, and to
delineate where the calculus genuinely extends those frameworks versus
where it specializes or restricts them.

\maplecourt{%
Maple Court's building model touches several of these frameworks
simultaneously.  The leak-humidity-mold Bayesian network is a special
case of both MLN and BN world models.  The ProbLog-style rule
``\texttt{0.3::mold\_risk(X) :- wall\_humidity(X, high)}'' encodes a
distributional fact that the WM calculus represents as extraction from
a state.  NARS-style evidence accumulation maps to the binary carrier.
Classical FOL is the degenerate case where all evidence is crisp.
Each bridge in this chapter will revisit this running example.}


%% ═════════════════════════════════════════════════════════════
\section{Markov Logic Networks}
\label{sec:bridge-mln}
%% ═════════════════════════════════════════════════════════════

Markov Logic Networks (MLNs)~\cite{richardson2006markov} combine
first-order logic with undirected graphical models.  An MLN is a set of
weighted first-order formulas $\{(w_i, F_i)\}_{i=1}^m$.  Given a finite
set of constants (the \emph{domain}), the MLN defines a ground Markov
network whose log-linear probability over possible worlds~$\omega$ is:
\[
  P(\omega) = \frac{1}{Z} \exp\!\Bigl(\sum_{i=1}^m w_i \cdot n_i(\omega)\Bigr)
\]
where $n_i(\omega)$ is the number of true groundings of formula~$F_i$ in
world~$\omega$, and $Z$ is the partition function.

\begin{definition}[MLN as world-model instance (\claimCore{})]
\label{def:mln-wm}
An MLN over domain~$D$ with formulas $\{(w_i, F_i)\}$ defines an
additive world model:
\begin{itemize}
\item \textbf{State:} a vector of grounding counts
  $W = (n_1, \ldots, n_m) \in \mathbb{R}_{\ge 0}^m$, where $n_i$
  records the weighted evidence for formula~$F_i$.
\item \textbf{Query:} a ground atom or formula $q$.
\item \textbf{Value type:} $V = \mathbb{R}$ (log-probabilities or
  sufficient statistics under the exponential family).
\item \textbf{Extraction:} $\mathrm{extract}(W, q) = \sum_{i : F_i
  \text{ relevant to } q} w_i \cdot n_i$ --- the total weighted
  evidence for~$q$ under the log-linear parameterization.
\item \textbf{Revision:} pointwise addition of grounding counts:
  $(n_1, \ldots, n_m) + (n_1', \ldots, n_m') =
  (n_1 + n_1', \ldots, n_m + n_m')$.
\end{itemize}
\end{definition}

\noindent
The extraction is additive because the log-linear model is additive in
sufficient statistics: adding evidence for formula~$F_i$ adds to the
weight of every world where $F_i$ holds.  The log-linear
parameterization is precisely what makes MLNs natural additive
world-model instances.

\begin{theorem}[MLN profile-hom correspondence (\claimCore{})]
\label{thm:mln-hom}
The MLN extraction map is an additive monoid homomorphism from
grounding-count states to log-probability profiles:
\[
  \widehat{\mathrm{extract}}_{\mathrm{MLN}} :
  \mathbb{R}_{\ge 0}^m \;\xrightarrow{\;+\;}\;
  (\mathrm{Query} \to \mathbb{R})
\]
This is a special case of the profile-hom equivalence
(Theorem~\ref{thm:universal-property}) with $V = \mathbb{R}$.
\end{theorem}

\paragraph{What the bridge preserves.}
The embedding preserves: additive revision, extraction linearity, and
the profile-hom structure.  An MLN state embeds faithfully into the
WM calculus with all five algebraic rules holding.

\paragraph{What the bridge does not preserve.}
The MLN bridge does \emph{not} preserve the partition function~$Z$
under revision.  Adding evidence changes the normalization constant,
so the map from WM states to normalized MLN probabilities is not
additive---it is a view projection, exactly as in
Chapter~\ref{ch:views}.  Normalized probability is a lossy view of
the underlying log-linear state.

\paragraph{The weight--evidence duality.}
MLN weights~$w_i$ play a dual role: they are both structural parameters
(defining the model) and evidence multipliers (scaling the contribution
of each formula).  In the WM calculus, these roles separate cleanly:
weights belong to the \emph{model class} (the $\Sigma$ side condition),
while counts belong to the \emph{state}.  This separation is why the
WM calculus can track provenance and support forgetting, while standard
MLNs cannot: in an MLN, retracting a single observation requires
recomputing the full partition function, because the observation's
contribution is entangled with the global normalization.

\maplecourt{%
The Maple Court leak-humidity-mold chain can be encoded as an MLN with
three weighted formulas:
$w_1 \cdot \mathrm{PipeLeak}(x)$,
$w_2 \cdot (\mathrm{PipeLeak}(x) \Rightarrow
\mathrm{WallHumidity}(x))$,
$w_3 \cdot (\mathrm{WallHumidity}(x) \Rightarrow
\mathrm{MoldRisk}(x))$.
The MLN grounding counts are $(n_1, n_2, n_3)$; the WM state is the
same vector.  Extraction for the query $\mathrm{MoldRisk}(3B)$ sums
the weighted contributions: $w_2 n_2 + w_3 n_3$.  The bridge maps
the WM compiled chain rule (Theorem~\ref{thm:chain-deduction}) to
the MLN's MAP inference under log-linear scoring.}


%% ═════════════════════════════════════════════════════════════
\section{ProbLog and Probabilistic Logic Programming}
\label{sec:bridge-problog}
%% ═════════════════════════════════════════════════════════════

ProbLog~\cite{DeRaedt2007ProbLog} annotates Prolog clauses with
probabilities: \texttt{0.3::leak(X) :- old\_pipe(X)} means that the
clause fires with probability~$0.3$.  The semantics is given by the
\emph{distribution semantics}~\cite{Sato1995}: each probabilistic fact
independently determines whether it is included in a logic program, and
the probability of a query is the sum over all \emph{explanations}
(sets of facts that entail the query).

\begin{definition}[ProbLog as world-model instance (\claimConditional{})]
\label{def:problog-wm}
A ground ProbLog program with $n$ probabilistic facts defines a
world-model instance:
\begin{itemize}
\item \textbf{State:} the vector of probabilistic-fact inclusion
  indicators, $(p_1, \ldots, p_n)$, enriched with observation counts
  for each fact.
\item \textbf{Query:} a ground atom $q$ derivable from the logic
  program.
\item \textbf{Value type:} $V = \mathbb{R}_{\ge 0}^2$ (binary evidence:
  positive count = number of worlds where $q$ holds, negative = where
  it does not).
\item \textbf{Extraction:} for each possible world~$\omega$
  (a subset of probabilistic facts), check whether $q$ is derivable
  under $\omega$'s logic program; accumulate binary evidence.
\end{itemize}
\end{definition}

\noindent
The side condition (\claimConditional{}) is that the ProbLog program is
\emph{ground and acyclic}: the distribution semantics requires a
well-founded derivation order.  Under this condition, the
extraction---summing over explanations---commutes with evidence
revision.

\paragraph{Distributional semantics vs.\ WM extraction.}
ProbLog's distribution semantics computes the probability of a query by
marginalizing over all possible worlds.  This is precisely what WM
extraction does for the Dirichlet-over-worlds instance
(Chapter~\ref{ch:world-models}, \S``A Worked Instance''):
\[
  \mathrm{extract}_{\mathrm{ProbLog}}(W, q) =
  \sum_{\omega \models q} P(\omega \mid W)
  = \mathrm{extract}_{\mathrm{WM}}(W, \mathrm{prop}(q))
\]
The WM extraction generalizes ProbLog's distribution semantics in two
directions:
\begin{enumerate}
\item ProbLog's probabilistic facts have \emph{fixed} probabilities;
  the WM calculus tracks evidence counts and computes probabilities as
  views.  This means WM states can be revised as new data arrives,
  while ProbLog programs must be re-parameterized.
\item ProbLog's queries are restricted to the Herbrand base of the
  logic program; the WM query type is generic and can include
  temporal, intensional, and deontic modalities
  (Chapters~\ref{ch:intensional}--\ref{ch:governance}).
\end{enumerate}

\begin{proposition}[ProbLog probability as WM view (\claimCore{})]
\label{prop:problog-view}
The probability assigned to a query~$q$ by a ProbLog program is the
strength view of the binary evidence extracted by the corresponding
WM instance:
\[
  P_{\mathrm{ProbLog}}(q) = \strength(\mathrm{extract}(W, q))
  = \frac{n_q^+}{n_q^+ + n_q^-}
\]
where $(n_q^+, n_q^-)$ is the binary evidence obtained by counting
worlds.
\end{proposition}

\paragraph{The ProbLog--BN connection.}
For ProbLog programs with a tree-structured dependency graph, the
distribution semantics reduces to exact BN inference.  The bridge
to BN world models (Chapter~\ref{ch:bn-models}) then applies
directly: compiled chain/fork/collider rewrites are available under
the d-separation side conditions induced by the ProbLog dependency
structure.

\maplecourt{%
The ProbLog rule
\texttt{0.3::mold\_risk(X) :- wall\_humidity(X, high)}
corresponds to a CPT entry in the Maple Court BN:
$P(\mathrm{MoldRisk} = \mathrm{true} \mid \mathrm{WallHumidity} =
\mathrm{high}) = 0.3$.  In the WM calculus, this CPT entry is not a
fixed parameter but a view of the evidence state
$(3, 7)$---three mold events in ten high-humidity observations.  As
more observations arrive, the evidence state is revised; the ProbLog
probability is the strength view of the current state.}


%% ═════════════════════════════════════════════════════════════
\section{Probabilistic Programming Languages}
\label{sec:bridge-ppl}
%% ═════════════════════════════════════════════════════════════

Probabilistic programming languages (PPLs) --- Church~\cite{Goodman2008Church},
Stan~\cite{Carpenter2017Stan}, WebPPL~\cite{Goodman2014WebPPL},
Pyro~\cite{Bingham2019Pyro}, and their kin --- elevate probabilistic
models to first-class programs.  A PPL program is a generative model:
it defines a joint distribution over all random choices made during
execution, and inference consists of conditioning that distribution on
observed data.  The dominant inference strategies are Markov Chain Monte
Carlo (MCMC) and variational inference (VI).

\paragraph{The PPL paradigm.}
In a PPL, the programmer writes a stochastic \emph{program}~$M$ that
samples from prior distributions and performs deterministic computation.
The posterior $P(\theta \mid \mathrm{obs})$ is the distribution over
latent variables~$\theta$ conditioned on observed data.  Inference
engines (NUTS, HMC, ADVI, etc.) explore this posterior by sampling or
by optimizing a variational approximation.  The program serves
simultaneously as model specification and as a generator of sample
execution traces.

\begin{definition}[PPL model as world-model instance (\claimConditional{})]
\label{def:ppl-wm}
A PPL model $M$ with latent space~$\Theta$ and observation
space~$\mathcal{X}$ defines a world-model instance:
\begin{itemize}
\item \textbf{State:} a sufficient-statistic summary of the
  accumulated data $\mathbf{x} = (x_1, \ldots, x_n)$, i.e.\
  $W = \mathrm{suffStats}(\mathbf{x}) \in V$.
\item \textbf{Query:} a function $q : \Theta \to \mathbb{R}$ (a
  posterior expectation or marginal density).
\item \textbf{Value type:} $V$ is the carrier appropriate for the
  model family (binary, Dirichlet, Normal-Gamma, etc.).
\item \textbf{Extraction:} $\mathrm{extract}(W, q) =
  \mathbb{E}_{\theta \sim P(\theta \mid W)}[q(\theta)]$, the
  posterior expectation of $q$ given the accumulated data.
\item \textbf{Revision:} $W_1 \hplus W_2 = $ componentwise addition
  of sufficient statistics (valid when observations are
  exchangeable --- de~Finetti~\cite{deFinetti1974}).
\end{itemize}
The side condition (\claimConditional{}) is exchangeability: if the
model assumes i.i.d.\ observations, revision is additive.  For
non-exchangeable models (time-series, structured priors), the overlap-aware
or trust-gated regimes apply.
\end{definition}

\paragraph{What the WM calculus adds over PPLs.}
Three structural gaps separate PPL-style inference from the WM calculus:

\begin{enumerate}
\item \textbf{Approximate vs.\ exact.}  PPL inference engines
  (NUTS, ADVI) approximate the posterior.  The WM calculus provides
  exact inference paths under explicit side conditions: the BN
  exactness theorems (Chapter~\ref{ch:bn-models}) certify that
  compiled rewrites agree with full extraction, with no sampling
  error.  For models with BN structure, the WM compiled path is
  both exact and fast; PPL samplers are approximate by design.

\item \textbf{No revision semantics.}  PPL programs are
  \emph{re-fit} when new data arrives: the model is re-run with a
  larger dataset.  The WM calculus provides an algebraic revision
  operation --- $W \hplus W'$ --- that updates the sufficient
  statistics without re-running the program.  This is online
  Bayesian updating; it requires no inference engine at
  revision time.  The sleep-consolidation theorem
  (Theorem~\ref{thm:sleep-consolidation}) guarantees that batch and
  sequential revision agree.

\item \textbf{Implicit side conditions.}  PPL models embed
  independence assumptions (conditional independence given latent
  variables, exchangeability) informally in the program structure.
  The WM calculus surfaces these as explicit $\Sigma$ predicates.
  The no-go theorem (Chapter~\ref{ch:no-go}) shows that ignoring
  side conditions leads to unsound inference; the $\Sigma$-discipline
  makes every compiled shortcut honest.
\end{enumerate}

\begin{proposition}[PPL posterior as WM strength view (\claimCore{})]
\label{prop:ppl-strength-view}
For a conjugate PPL model (Beta-Binomial, Dirichlet-Multinomial, or
Normal-Gamma), the posterior mean of any query~$q$ is the WM strength
view of the corresponding evidence:
\[
  \mathbb{E}[q(\theta) \mid \mathbf{x}] =
  \strength(\mathrm{extract}(W, q))
\]
where $W = \mathrm{suffStats}(\mathbf{x})$.  The full posterior
distribution is the view returned by \emph{distributional extraction}
(Chapter~\ref{ch:error-transport}), not a scalar.
\end{proposition}

\paragraph{PPLBench and the nine-lane validation.}
The PPLBench comparison (Chapter~\ref{ch:pplbench}) runs the WM
calculus head-to-head against nine PPL implementations (Stan, Pyro,
Turing.jl, etc.) on a common benchmark suite.  The WM calculus
achieves exact posterior agreement with Stan on conjugate models
(where Stan's NUTS converges to the exact posterior), while using
only sufficient-statistic arithmetic and zero sampling.  For
non-conjugate models, the WM distributional path
(Chapter~\ref{ch:error-transport}) provides certified approximation
bounds; PPL samplers provide convergence guarantees in the limit
but no per-query error certificates.

\paragraph{The PPL--BN boundary.}
A PPL model that compiles to a BN (e.g., a Pyro model with an acyclic
data-generating process and conjugate priors) falls under the exact BN
inference regime of Chapter~\ref{ch:bn-models}.  In this case:
\begin{itemize}
\item The PPL's MCMC is unnecessary: the BN query bridge
  (Theorem~\ref{thm:chain-deduction}) gives the exact posterior.
\item The WM compiled rules (chain, fork, collider) replace NUTS
  as the inference engine.
\item Revision is evidence addition; re-fitting is unnecessary.
\end{itemize}
This boundary delineates where the WM calculus dominates PPL inference
outright (BN-structured conjugate models) versus where PPLs remain
competitive (non-conjugate, non-BN models requiring VI or MCMC).

\maplecourt{%
The Maple Court humidity model can be expressed as a Stan program:
a Normal prior on the mean humidity, updated by observations from
the hallway sensor.  Stan's NUTS would run $\sim\!10^4$ steps to
estimate the posterior mean.  The WM calculus accumulates sufficient
statistics $(n, \sum x_i, \sum x_i^2)$ in the Normal-Gamma carrier
and computes the posterior mean analytically:
$\mu_n = (\kappa_0 \mu_0 + \sum x_i) / (\kappa_0 + n)$.
Both approaches agree on the posterior mean; the WM calculus
provides the exact answer in $O(1)$, while Stan provides an
approximation in $O(n_{\mathrm{samples}})$.  For the
non-conjugate question ``what is the probability that \emph{both}
pipe-leak evidence \emph{and} humidity jointly exceed their
respective thresholds simultaneously?'' --- a query that involves
the joint distribution across two carriers --- the WM distributional
path applies, while Stan would use a dedicated multi-variate model.}


%% ═════════════════════════════════════════════════════════════
\section{PLN v0.9 Subsumption}
\label{sec:bridge-pln}
%% ═════════════════════════════════════════════════════════════

The world-model calculus grew from the formalization of Probabilistic
Logic Networks~\cite{goertzel2009pln}.  This section makes the
subsumption precise: every classical PLN rule is a compiled WM
instance, exact under explicit side conditions.

\paragraph{The classical rule family.}
PLN v0.9 provides seven core inference rules: deduction, induction,
abduction, modus ponens, revision, negation, and inversion.  Each rule
takes truth values (strength, confidence pairs) as inputs and produces
a truth value as output.  The rules are stated as algebraic formulas
over strength and confidence.

The world-model calculus subsumes this family in the following precise
sense:

\begin{theorem}[PLN subsumption (\claimConditional{})]
\label{thm:pln-subsumption}
Every PLN v0.9 rule is a $\Sigma$-guarded compiled rewrite
(Chapter~\ref{ch:sigma-rewrites}) of the additive world-model
calculus.  Specifically:
\begin{enumerate}
\item Each rule's strength formula is the strength \emph{view} of a
  full posterior extraction under stated side conditions
  (screening-off, positivity, stamp-disjointness).
\item Each rule's confidence formula is the confidence view of the
  same extraction.
\item Each rule is exact (the compiled rewrite agrees with full
  extraction) when its $\Sigma$ holds; when $\Sigma$ fails, the
  rule is not licensed.
\end{enumerate}
\end{theorem}

\noindent
The proof is distributed across Chapters~\ref{ch:sigma-rewrites}
and~\ref{ch:bn-models}: the deduction formula is proved exact under
chain-BN screening-off (Theorem~\ref{thm:chain-deduction}); induction
and abduction reduce to Bayes-inversion followed by deduction; modus
ponens, revision, negation, and inversion each have their own
$\Sigma$-proofs.

\paragraph{What the calculus adds.}
The WM calculus adds five capabilities that PLN v0.9 lacks:
\begin{enumerate}
\item \textbf{Explicit side conditions.}  PLN v0.9 states rules as
  universally valid formulas; the WM calculus states the conditions
  under which each rule is exact.
\item \textbf{Distributional inference.}  PLN v0.9 propagates scalar
  (strength, confidence) summaries; the WM calculus propagates full
  posterior distributions (Chapter~\ref{ch:error-transport}), with
  five proved limitation theorems quantifying the cost of scalar
  chaining.
\item \textbf{Provenance and forgetting.}  PLN v0.9 carries evidence
  stamps but cannot formally retract evidence; the WM calculus
  provides scoped forgetting with conservation guarantees
  (Chapter~\ref{ch:state-ops}).
\item \textbf{Multiple evidence carriers.}  PLN v0.9 uses binary
  evidence exclusively; the WM calculus supports Dirichlet,
  Normal-Gamma, and generic carriers (Chapter~\ref{ch:evidence}).
\item \textbf{Typed rewrite verification.}  PLN v0.9 justifies rules
  by informal argument; the WM calculus provides kernel-checked Lean
  proofs for every rule in the catalog.
\end{enumerate}

\paragraph{What PLN v0.9 provides that the calculus inherits.}
The original PLN design contributes the \emph{conceptual vocabulary}
that the WM calculus formalizes: evidence-based reasoning, the
strength-confidence decomposition, the idea that inference rules are
derived rather than primitive, and the revision-as-evidence-combination
pattern.  The WM calculus does not replace this vocabulary---it gives
it a precise typed semantics.

\begin{remark}[PLN v0.9 examples as WM conformance tests]
The three flagship PLN v0.9 examples---DeductionRevision,
RavenInduction, and FlyingRaven---are fully verified as compiled
BN rewrites with provenance tracking and 20~kernel-checked
conformance tests (Chapter~\ref{ch:bn-models}, \S``PLN v0.9
Examples'').
\end{remark}


%% ═════════════════════════════════════════════════════════════
\section{NARS Bridge}
\label{sec:bridge-nars}
%% ═════════════════════════════════════════════════════════════

The Non-Axiomatic Reasoning System (NARS) of Pei Wang uses a
truth-value system based on two components: \emph{frequency}
$f = w^+ / (w^+ + w^-)$ and \emph{confidence}
$c = (w^+ + w^-) / (w^+ + w^- + k)$, where $w^+$ and $w^-$ are
positive and negative evidence counts and $k$ is a system parameter
(typically $k = 1$).

\begin{definition}[NARS truth value as WM binary evidence (\claimCore{})]
\label{def:nars-bridge}
The NARS truth value $(f, c)$ with parameter~$k$ corresponds to the
WM binary evidence $(n^+, n^-)$ via:
\[
  n^+ = \frac{f \cdot c \cdot k}{1 - c}, \qquad
  n^- = \frac{(1-f) \cdot c \cdot k}{1 - c}
\]
Conversely, $f = n^+ / (n^+ + n^-)$ and
$c = (n^+ + n^-) / (n^+ + n^- + k)$.
\end{definition}

\noindent
The correspondence is exact: the NARS frequency is the WM strength
view, and the NARS confidence is the WM confidence view with
denominator offset~$k$ instead of~$2$ (the WM Laplace smoothing
uses $k = 2$; NARS typically uses $k = 1$).

\begin{theorem}[NARS revision as WM revision (\claimCore{})]
\label{thm:nars-revision}
The NARS revision rule---which combines two truth values
$(f_1, c_1)$ and $(f_2, c_2)$ into a weighted average---is
equivalent to additive evidence revision in the WM calculus:
\[
  (n_1^+, n_1^-) \hplus (n_2^+, n_2^-) = (n_1^+ + n_2^+, n_1^- + n_2^-)
\]
The NARS revision formula
$f_{\mathrm{rev}} = (f_1 c_1 + f_2 c_2) / (c_1 + c_2)$
is the strength view of the summed evidence, provided the evidence
sets are disjoint (the stamp-disjointness condition of
Chapter~\ref{ch:sigma-rewrites}).
\end{theorem}

\begin{proof}[Proof sketch]
Converting NARS truth values to evidence counts:
$w_i = c_i k / (1 - c_i)$, $w_i^+ = f_i w_i$, $w_i^- = (1-f_i) w_i$.
Summing and converting back:
$f_{\mathrm{rev}} = (w_1^+ + w_2^+) / (w_1 + w_2)
= (f_1 w_1 + f_2 w_2) / (w_1 + w_2)$.
Since $w_i$ is monotonically related to $c_i$ (for fixed~$k$), the
NARS weighted-average formula is the strength view of additive
evidence combination.
\end{proof}

\paragraph{Where NARS and WM diverge.}
Three points of divergence:
\begin{enumerate}
\item \textbf{Inference rules.}  NARS deduction uses the formula
  $f_{AC} = f_{AB} \cdot f_{BC}$ with a confidence adjustment.  The
  WM deduction formula (Definition~\ref{def:deduction-strength}) is
  more complex because it accounts for the base rate:
  $s_{AC} = s_{AB} s_{BC} + (1 - s_{AB})(s_C - s_B s_{BC})/(1 - s_B)$.
  The NARS formula is the special case where the base rate terms
  vanish (i.e., $s_C = s_B \cdot s_{BC}$).
\item \textbf{Side conditions.}  NARS rules fire unconditionally;
  the WM calculus requires explicit $\Sigma$ for each rule.  The
  no-go theorem (Chapter~\ref{ch:no-go}) shows that unconditional
  firing is unsound in general.
\item \textbf{Temporal handling.}  NARS has native temporal operators
  (predictive and retrospective implications); the WM calculus
  handles temporality via the temporal quantale
  (Chapter~\ref{ch:temporal}).  Both approaches use sequential
  composition, but the WM quantale structure provides free
  transitivity theorems (Theorem~\ref{thm:temporal-transitivity})
  that NARS obtains by separate inference rules.
\end{enumerate}

\maplecourt{%
Suppose a NARS agent in Maple Court observes five leak alarms and
two no-alarm readings for Apartment~3B's pipe sensor.  The NARS
truth value is $f = 5/7 \approx 0.71$, $c = 7/8 = 0.875$ (with
$k = 1$).  The WM binary evidence is $(5, 2)$.  Both systems
agree on the evidence; they diverge when chaining to
$\mathrm{MoldRisk}$: NARS applies $f_{AC} = f_{AB} \cdot f_{BC}$
unconditionally, while the WM calculus requires screening-off and
uses the full deduction formula.  For this particular BN structure,
the WM formula reduces to the NARS formula only when the mold base
rate satisfies $s_C = s_B \cdot s_{BC}$---a condition that the
$\Sigma$-discipline makes explicit.}


%% ═════════════════════════════════════════════════════════════
\section{Classical FOL and Set Theory as Degenerate WM}
\label{sec:bridge-fol}
%% ═════════════════════════════════════════════════════════════

Classical first-order logic and set theory are world-model instances
in a degenerate but illuminating way.

\begin{definition}[Boolean world model (\claimCore{})]
\label{def:boolean-wm}
A \emph{Boolean world model} is an additive world model with value
type $V = \mathbb{B} = \{0, 1\}$ under Boolean OR as addition:
\begin{itemize}
\item $\mathrm{State} = \mathcal{P}(\Omega)$ (a set of possible
  worlds, or equivalently a first-order theory).
\item $\mathrm{Query} = $ first-order sentences over the signature.
\item $\mathrm{extract}(W, q) = 1$ if $W \models q$, and $0$
  otherwise.
\item $\mathrm{revise}(W_1, W_2) = W_1 \cap W_2$ (intersection of
  model sets, i.e., conjunction of theories).
\end{itemize}
\end{definition}

\noindent
In this instance, the ``additive'' structure degenerates: Boolean OR
is idempotent ($1 + 1 = 1$), so the additive monoid collapses to a
bounded semilattice.  Extraction is the characteristic function of
the satisfaction relation.

\begin{theorem}[FOL as degenerate additive WM (\claimCore{})]
\label{thm:fol-wm}
The Boolean world model satisfies the additive world-model axioms
with $V = (\mathbb{B}, \vee, 0)$:
\begin{enumerate}
\item $\mathrm{extract}(W_1 \cap W_2, q) =
  \mathrm{extract}(W_1, q) \vee \mathrm{extract}(W_2, q)$ when
  restricted to queries that are monotone in the model set.
\item $\mathrm{extract}(\Omega, q) = 0$ for non-tautological queries
  (the empty state knows nothing).
\end{enumerate}
For classical (monotone) entailment, all five algebraic rules hold.
\end{theorem}

\paragraph{Set-theoretic operations as WM operations.}
The standard set-theoretic operations have WM counterparts:
\begin{center}
\begin{tabular}{ll}
\toprule
\textbf{Set theory / FOL} & \textbf{WM calculus} \\
\midrule
$W \models \varphi$ (entailment) & $\mathrm{extract}(W, \varphi) = 1$ \\
$\mathrm{Th}(\mathcal{M})$ (theory of a model) & state \\
$\mathrm{Mod}(T)$ (models of a theory) & $\{W : \mathrm{extract}(W,q)=1\}$ \\
$T_1 \cup T_2$ (theory union) & $\mathrm{revise}(W_1, W_2)$ \\
$\emptyset$ (empty theory) & $\mathrm{empty}$ \\
\bottomrule
\end{tabular}
\end{center}

\paragraph{Why this bridge matters.}
The degenerate bridge shows that the WM calculus \emph{contains}
classical logic as a special case, but with a crucial loss: the
Boolean carrier discards all uncertainty information.  There is no
strength, no confidence, no evidence count---only yes/no.
The information ordering collapses to a total order ($0 < 1$), and
the lattice structure of the evidence carriers (Figure~\ref{fig:evidence-lattice})
reduces to a two-element chain.

This is precisely the limitation that motivated the calculus:
classical logic can express what is known but not how well it is
known.  The WM calculus replaces the Boolean carrier with richer
carriers that track both \emph{what} and \emph{how much}.

\begin{proposition}[Crisp extraction (\claimCore{})]
\label{prop:crisp}
For any additive world model with carrier~$V$, the ``crisp
extraction'' obtained by composing extraction with the threshold
view $\mathbb{1}_{s > 0.5}$ recovers Boolean entailment:
\[
  \mathrm{crisp}(W, q) = \begin{cases}
    1 & \text{if } \strength(\mathrm{extract}(W, q)) > 0.5 \\
    0 & \text{otherwise}
  \end{cases}
\]
This is a view projection (Chapter~\ref{ch:views}): it discards
all information except the sign of the evidence balance.
\end{proposition}


%% ═════════════════════════════════════════════════════════════
\section{Categorical Probability Bridge}
\label{sec:bridge-categorical}
%% ═════════════════════════════════════════════════════════════

The categorical probability program of Fritz, Rischel, and
collaborators (see Fritz 2020, \emph{A synthetic approach to Markov
kernels}) provides a setting in which probability theory is developed
abstractly using \emph{Markov categories}: symmetric monoidal
categories in which every object is equipped with a copy and discard
morphism satisfying naturality conditions.

\paragraph{Markov categories and the WM calculus.}
A Markov category $\mathbf{C}$ has:
\begin{itemize}
\item Objects: measurable spaces (or their abstract analogues).
\item Morphisms: Markov kernels $X \to Y$ (stochastic maps).
\item Monoidal product: independent combination of spaces.
\item Copy morphism: $\mathrm{copy}_X : X \to X \otimes X$ (duplication).
\item Discard morphism: $\mathrm{del}_X : X \to I$ (marginalization).
\end{itemize}

\begin{definition}[WM extraction as Markov-category morphism (\claimCore{})]
\label{def:markov-cat-bridge}
In a Markov category, the WM extraction map
$\mathrm{extract} : \mathrm{State} \times \mathrm{Query} \to V$
is a \emph{deterministic} morphism (a morphism that commutes with
copy): querying the same state twice for the same question yields
the same answer.  The additivity law
$\mathrm{extract}(W_1 + W_2, q) = \mathrm{extract}(W_1, q) +
\mathrm{extract}(W_2, q)$ expresses that extraction is a
\emph{monoid homomorphism in the Markov-category sense}: it
preserves the monoidal structure of state combination.
\end{definition}

\begin{theorem}[Profile-hom in Markov categories (\claimCore{})]
\label{thm:markov-cat-hom}
The profile-hom equivalence (Theorem~\ref{thm:universal-property})
has a categorical restatement: an additive world model is an
\emph{algebra for the commutative-monoid monad} in the category of
additive monoids, with the extraction being the algebra map.  In a
Markov category, this becomes: an additive world model is a
deterministic morphism from the state monoid to the function monoid
$\mathrm{Query} \to V$, and this morphism commutes with copy
(determinism) and discard (zero extraction).
\end{theorem}

\paragraph{The Giry monad connection.}
The Giry monad $\mathcal{G}$ sends a measurable space~$X$ to the
space of probability measures on~$X$.  The Kleisli category of~$\mathcal{G}$
is a Markov category.  In this setting:
\begin{itemize}
\item A WM state is a probability measure on the space of possible
  worlds: $W \in \mathcal{G}(\Omega)$.
\item Revision is the convolution (mixture) of measures---the Giry
  monad's multiplication applied to the sum of evidence.
\item The Giry bind implements Kyburg flattening
  (Chapter~\ref{ch:world-models}, \S``Giry monad and Markov
  kernels'').
\end{itemize}
The formalization proves that this flattening satisfies all three
monad laws, establishing the probability-theoretic foundation
categorically.

\paragraph{Copy-discard categories and evidence.}
The evidence carriers of Chapter~\ref{ch:evidence} live in a
\emph{copy-discard category}: the copy morphism on binary evidence
$(n^+, n^-)$ duplicates the counts (making two copies of the same
evidence), and the discard morphism sends all evidence to $(0, 0)$
(the empty state).  The quantale tensor
(Theorem~\ref{thm:quantale-trans}) is a morphism in this category:
it composes evidence along inference chains while preserving the
copy-discard structure.

\begin{remark}[Categorical coherence]
The categorical probability bridge places the WM calculus in the
landscape of 2020s categorical probability.  The profile-hom
equivalence is the algebraic shadow of a Markov-category structure;
the Giry monad provides the measure-theoretic foundation; the
copy-discard laws ensure that evidence behaves well under duplication
and marginalization.  The open conjecture
(\S\ref{sec:gslt-open-program}) that every GSLT induces a world
model can be restated categorically: every graph-structured lambda
theory with parallel composition defines a commutative-monoid
object in a suitable Markov category.
\end{remark}


%% ═════════════════════════════════════════════════════════════
\section{Process Calculus Bridge}
\label{sec:bridge-process}
%% ═════════════════════════════════════════════════════════════

Meredith's $\rho$-calculus (the reflective higher-order process
calculus) provides a computational foundation in which processes are
first-class citizens: they can be sent as messages, composed in
parallel, and inspected structurally.  The connection to the
world-model calculus runs through the GSLT universality theorem
and the amplitude carrier.

\paragraph{The $\rho$-calculus in brief.}
The $\rho$-calculus~\cite{MeredithRadestock2005} extends the
$\pi$-calculus with reflection: a process $P$ can be ``reified'' as
a name $\ulcorner P \urcorner$ and sent over a channel, and
conversely a name can be ``reflected'' back into a running process.
The operational semantics has three constructors:
\begin{enumerate}
\item $0$ --- the stopped process (the empty state).
\item $P \mid Q$ --- parallel composition (revision).
\item $\mathrm{for}(y \leftarrow x)\{P\}$ --- input-guarded
  continuation (extraction awaiting a query).
\end{enumerate}

\begin{definition}[$\rho$-calculus as world-model instance (\claimCore{})]
\label{def:rho-wm}
A $\rho$-calculus process defines a world-model instance:
\begin{itemize}
\item $\mathrm{State} = \mathrm{Process}$ (the set of
  $\rho$-calculus processes modulo structural congruence).
\item $\mathrm{Query} = \mathrm{Name}$ (channel names, which are
  quoted processes).
\item $V$ = a suitable observation type (barbs, or more generally
  the values observed on channels).
\item $\mathrm{revise}(P, Q) = P \mid Q$ (parallel composition).
\item $\mathrm{empty} = 0$ (the stopped process).
\item $\mathrm{extract}(P, x) = $ the observable behavior of $P$
  on channel~$x$ (the barb $P \downarrow_x$).
\end{itemize}
\end{definition}

\noindent
Parallel composition is commutative and associative (by structural
congruence), so the state space is a commutative monoid.  The
stopped process is the unit.  The extraction (barb observation) is
a form of query answering: ``does process $P$ exhibit behavior on
channel~$x$?''

\paragraph{The amplitude regime.}
When $V = \mathbb{C}$ (complex amplitudes), the WM calculus enters
what Meredith calls the \emph{quantum amplitude regime}.  In this
regime:
\begin{itemize}
\item Evidence values are complex amplitudes rather than real counts.
\item Revision is addition of amplitudes (superposition).
\item The Born rule $|z|^2$ is a view projection from amplitudes to
  probabilities---exactly the lossy view pattern of
  Chapter~\ref{ch:views}.
\item Interference effects arise naturally: two evidence
  contributions can cancel ($z_1 + z_2 = 0$) or reinforce
  ($|z_1 + z_2|^2 > |z_1|^2 + |z_2|^2$).
\end{itemize}

\begin{proposition}[Amplitude extraction is additive (\claimCore{})]
\label{prop:amplitude-additive}
With $V = \mathbb{C}$ and revision as complex addition, the
extraction law holds:
\[
  \mathrm{extract}(P \mid Q, x) =
  \mathrm{extract}(P, x) + \mathrm{extract}(Q, x)
\]
The Born-rule view $|\cdot|^2$ is \emph{not} additive:
$|z_1 + z_2|^2 \ne |z_1|^2 + |z_2|^2$ in general.  This is another
instance of the views-are-not-additive principle
(Proposition~\ref{prop:overview-views-not-semantic-core}).
\end{proposition}

\paragraph{GSLT universality and process algebra.}
The GSLT universality theorem
(Theorem~\ref{thm:universalEnsembleWorldModel} in
Chapter~\ref{ch:typed}) says: for any type~$T$, multiset ensembles
of~$T$ give an additive world model.  A $\rho$-calculus process is a
term of type $\mathrm{Process}$; the multiset of its parallel
components is an ensemble; and decidable predicates on processes
(``does this process contain a send on channel~$x$?'') are queries.
The GSLT theorem therefore gives every $\rho$-calculus system a
canonical world model.

\begin{lstlisting}[language={}]
-- Rho-calculus process as WM state
instance : AdditiveWorldModel Process Name Observation where
  revise  := Process.par        -- parallel composition
  empty   := Process.stop       -- stopped process
  extract := Process.barb       -- observable behavior
  -- Laws: par_comm, par_assoc, par_zero, barb_par_add
\end{lstlisting}

\paragraph{The spatial-behavioral bridge.}
Meredith's spatial-behavioral types~\cite{meredith2026ccc} classify
processes by the spatial structure of their communication patterns.
In WM terms, spatial-behavioral types are \emph{query families}:
each type specifies a class of observations (barbs) that the process
may exhibit.  The type-checking judgment ``process $P$ has
spatial-behavioral type $\tau$'' becomes a WM extraction:
$\mathrm{extract}(P, \tau) \ne 0$ iff $P$ exhibits the behavior
classified by~$\tau$.

The OSLF modal operators $\Diamond$ and $\Box$
(Chapter~\ref{ch:typed}) are the WM calculus's version of
Hennessy-Milner modalities~\cite{HennessyMilner1985}: $\Diamond\varphi$
asks whether some reduction leads to a state satisfying~$\varphi$
(forward reachability), and $\Box\varphi$ asks whether every
predecessor satisfies~$\varphi$ (backward invariance).  The Galois
connection $\Diamond \dashv \Box$ is the process-algebraic
completeness theorem recast in WM terms.

\maplecourt{%
Imagine Maple Court's building automation as a concurrent system:
each sensor is a process emitting readings on its channel, and the
building controller is a process listening on all sensor channels.
The parallel composition $\mathrm{Sensor}_1 \mid \mathrm{Sensor}_2
\mid \cdots \mid \mathrm{Controller}$ is the WM state.  Querying
``is there a leak alarm on channel \texttt{pipe\_3B}?'' is a barb
observation.  The amplitude regime would model quantum sensors
(e.g., SQUID-based leak detectors) where interference between
sensor signals is physically meaningful.}


%% ═════════════════════════════════════════════════════════════
\section{What the Calculus Covers and Where It Stops}
\label{sec:bridge-scope}
%% ═════════════════════════════════════════════════════════════

The eight bridges above map out the calculus's neighborhood.  This
section summarizes the coverage and identifies the boundaries---a
preview of the full Limitations chapter (Chapter~\ref{ch:limitations}).

\paragraph{What the calculus covers.}
The world-model calculus provides a unified framework for:

\begin{enumerate}
\item \textbf{Evidence-based uncertain reasoning} with explicit
  carriers, revisable states, and compositional extraction
  (Chapters~\ref{ch:world-models}--\ref{ch:evidence}).
\item \textbf{Compiled inference} with explicit side conditions and
  kernel-checked correctness proofs, subsuming the PLN rule family
  and exact BN inference (Chapters~\ref{ch:sigma-rewrites}--\ref{ch:bn-models}).
\item \textbf{Distributional inference} that avoids the variance
  accumulation, sensitivity amplification, and coverage collapse of
  scalar chaining (Chapter~\ref{ch:error-transport}).
\item \textbf{State management}: forgetting, provenance,
  explanation, and conservation
  (Chapter~\ref{ch:state-ops}).
\item \textbf{Extensional and intensional inheritance} with a
  typed three-channel model refuting the classical scalar mixture
  (Chapter~\ref{ch:intensional}).
\item \textbf{Temporal and causal inference} via the temporal
  quantale, with free transitivity theorems
  (Chapter~\ref{ch:temporal}).
\item \textbf{Normative reasoning} with graded deontic evidence and
  a temporal-deontic bridge (Chapter~\ref{ch:governance}).
\item \textbf{Neighboring frameworks}: MLNs, ProbLog, NARS,
  classical FOL, categorical probability, and process calculi, all
  as instances or projections of the same calculus (this chapter).
\end{enumerate}

\paragraph{Where the calculus stops.}
Five boundaries delimit the current scope:

\begin{description}
\item[Non-additive regimes.] The overlap-aware, trust-gated, and
  coalition revision regimes are typed but not developed in full
  generality.  No profile-hom equivalence is known outside the
  additive regime (Chapter~\ref{ch:world-models}).

\item[Measure-theoretic grounding.] The Giry monad connection is
  stated but the terminal
  $\code{WorldModel (Measure\;\Omega)\;Query}$ instance remains
  future work.  The current formalization works with finitary
  carriers and discrete state spaces.

\item[Higher-order completeness.] The calculus handles quantified and
  higher-order queries as instances of the generic Query type, but
  no completeness theorem is proved for the full higher-order
  fragment.

\item[Computational complexity.] The calculus says \emph{what} the
  correct answer is but does not provide complexity bounds for
  computing it.  Exact BN inference is \#P-hard in general; the
  WM calculus licenses compiled shortcuts but does not analyze their
  computational cost.

\item[The probability hypercube.] The full classification of theories
  by which combination of merge-law and carrier properties they
  satisfy (commutativity, distributivity, precision, additivity) is
  a work in progress.  The 29-file Lean formalization in
  \codepath{Mettapedia/ProbabilityTheory/Hypercube/} maps the axes
  but does not yet prove the complete classification theorem.
\end{description}

\begin{figure}[htbp]
\centering
\begin{tikzpicture}[
  >=Latex,
  every node/.style={font=\footnotesize},
  framework/.style={draw, rounded corners=4pt, minimum width=28mm,
              minimum height=10mm, align=center, fill=blue!6},
  wmbox/.style={draw, very thick, rounded corners=6pt, minimum width=36mm,
              minimum height=12mm, align=center, fill=green!8},
  boundary/.style={draw, dashed, rounded corners=4pt, minimum width=28mm,
              minimum height=10mm, align=center, fill=orange!8},
]
  % Center: WM calculus
  \node[wmbox] (wm) at (0, 0) {WM Calculus\\[-2pt]{\scriptsize additive regime}};

  % Bridges (exact instances)
  \node[framework] (mln) at (-5, 2.5) {MLN};
  \node[framework] (problog) at (-2.5, 4.0) {ProbLog};
  \node[framework] (ppl) at (0, 4.8) {PPLs};
  \node[framework] (pln) at (2.5, 4.0) {PLN v0.9};
  \node[framework] (nars) at (5, 2.5) {NARS};
  \node[framework] (fol) at (-5, -2.5) {FOL / Set};
  \node[framework] (markov) at (0, -3.5) {Markov Cat.};
  \node[framework] (rho) at (5, -2.5) {$\rho$-calculus};

  % Arrows
  \draw[->, thick] (mln) -- node[above left, font=\tiny] {log-linear} (wm);
  \draw[->, thick] (problog) -- node[left, font=\tiny] {dist.\ sem.} (wm);
  \draw[->, thick] (ppl) -- node[right, font=\tiny] {suff.\ stats.} (wm);
  \draw[->, thick] (pln) -- node[right, font=\tiny] {$\Sigma$-compile} (wm);
  \draw[->, thick] (nars) -- node[above right, font=\tiny] {evidence} (wm);
  \draw[->, thick] (fol) -- node[below left, font=\tiny] {Boolean} (wm);
  \draw[->, thick] (markov) -- node[right, font=\tiny] {Giry monad} (wm);
  \draw[->, thick] (rho) -- node[below right, font=\tiny] {GSLT} (wm);

  % Boundaries
  \node[boundary] (nonadd) at (-6, 0) {non-additive\\[-2pt]{\scriptsize (frontier)}};
  \node[boundary] (measure) at (6, 0) {measure-theoretic\\[-2pt]{\scriptsize (open)}};

  \draw[->, dashed, orange!70!black] (nonadd) -- (wm);
  \draw[->, dashed, orange!70!black] (measure) -- (wm);
\end{tikzpicture}
\caption{The world-model calculus and its bridges.  Solid arrows:
exact embeddings (this chapter).  Dashed arrows: frontier directions
not yet fully developed.  PPLs (Church, Stan, Pyro, etc.) embed via
sufficient-statistic carriers under exchangeability.}
\label{fig:bridges-landscape}
\end{figure}

\noindent
The central design principle is visible in the diagram: the WM
calculus is not one more framework competing with MLNs, ProbLog,
NARS, or process calculi.  It is the \emph{typed evidence layer}
that each of these frameworks implicitly uses.  Making that layer
explicit---with revisable states, compositional extraction, explicit
side conditions, and kernel-checked proofs---is what the calculus
contributes.  Where it stops, it stops honestly: the boundaries are
stated as theorems (the no-go results) or as open programs (the
probability hypercube, GSLT universality, measure-theoretic
grounding), not as hand-waving.

\maplecourt{%
Maple Court's building model can be expressed in any of the eight
neighboring frameworks:
as an MLN with weighted formulas,
as a ProbLog program with probabilistic clauses,
as a Stan/Pyro PPL with a generative model,
as a PLN v0.9 rule application,
as a NARS evidence-accumulation system,
as a classical FOL theory (ignoring uncertainty),
as a Markov-category morphism (the categorical view), or
as a concurrent process system (the process-calculus view).
Each bridge in this chapter shows how the WM calculus unifies these
perspectives: the underlying state, evidence, extraction, and
revision are the same; only the surface language and the
computational strategy change.  The $\Sigma$-discipline ensures
that whichever view is used, the correctness conditions are
explicit.}
